% Part: computability
% Chapter: recursive-functions
% Section: trees

\documentclass[../../../include/open-logic-section]{subfiles}

\begin{document}

\olfileid[pt]{cmp}{rec}{tre}
\olsection{Árvores}

Às vezes é útil representar árvores como números naturais, assim como
podemos representar sequências por números, e suas propriedades e
operações por relações e funções recursivas primitivas sobre seus
códigos. Usaremos sequências e seus códigos para isso. Uma árvore pode
ser um único nó (possivelmente com um rótulo) ou um nó (possivelmente
com um rótulo) ligado a várias subárvores. O nó é chamado de
\emph{raiz} da árvore, e as subárvores às quais ele está ligado são
suas \emph{subárvores imediatas}.

Codificamos árvores recursivamente como uma sequência
$\tuple{k, d_1, \dots, d_k}$, em que $k$ é o número de subárvores
imediatas e $d_1$, \dots,~$d_k$ são os códigos das subárvores imediatas.
Se os nós tiverem rótulos, estes podem ser incluídos depois das
subárvores imediatas. Assim, uma árvore formada apenas por um único nó
com rótulo~$l$ seria codificada por $\tuple{0,l}$, e uma árvore formada
por uma raiz (rotulada~$l_1$) ligada a dois nós isolados (rotulados
$l_2$, $l_3$) seria codificada por $\tuple{2,
  \tuple{0, l_2}, \tuple{0, l_3}, l_1}$.

\begin{prop}
  \ollabel{prop:subtreeseq}
  A função $\fn{SubtreeSeq}(t)$, que devolve o código de uma
  sequência cujos elementos são os códigos de todas as subárvores da
  árvore de código~$t$, é recursiva primitiva.
\end{prop}

\begin{proof}
  Observe primeiro que $\fn{ISubtrees}(t) = \fn{subseq}(t, 1, (t)_0)$
  é recursiva primitiva e devolve os códigos das subárvores imediatas
  de uma árvore~$t$. Agora podemos definir uma função auxiliar
  $\fn{hSubtreeSeq}(t,n)$ que calcula a sequência de todas as subárvores
  que estão a $n$ nós de distância da raiz. A sequência das subárvores
  de~$t$ que estão a $0$ nós de distância da raiz---em outras palavras,
  que começa na raiz de~$t$---é a sequência formada apenas por~$t$.
  Para obter uma sequência de todas as subárvores de nível~$n+1$ de
  $t$, concatenamos as subárvores de nível $n$ com uma sequência formada
  por todas as subárvores imediatas das subárvores de nível $n$. Para
  obter uma lista de todas elas, observe que, se $f(x)$ é uma função
  recursiva primitiva que devolve códigos de sequências, então
  $g_f(s, k) = f((s)_0) \concat \dots \concat f((s)_k)$ também é
  recursiva primitiva:
    \begin{align*}
      g(s, 0) & = f((s)_0)\\
      g(s, k+1) & = g(s, k) \concat f((s)_{k+1})
    \end{align*}
    Por exemplo, se $s$ é uma sequência de árvores, então
    $h(s) = g_{\fn{ISubtrees}}(s, \len{s})$ fornece a sequência das
    subárvores imediatas dos elementos de~$s$. Podemos usá-la para
    definir $\fn{hSubtreeSeq}$ por
    \begin{align*}
      \fn{hSubtreeSeq}(t, 0) & = \tuple{t} \\
      \fn{hSubtreeSeq}(t, n+1) & = \fn{hSubtreeSeq}(t, n) \concat
      h(\fn{hSubtreeSeq}(t, n)).
    \end{align*}
    O nível máximo das subárvores em uma árvore codificada por~$t$,
    isto é, a distância máxima entre a raiz e um nó folha, é limitado
    pelo código~$t$. Logo, uma sequência dos códigos de todas as
    subárvores da árvore codificada por~$t$ é dada por
    $\fn{hSubtreeSeq}(t, t)$.
\end{proof}

\begin{prob}
  A definição de $\fn{hSubtreeSeq}$ na demonstração de
  \olref[cmp][rec][tre]{prop:subtreeseq} inclui, em geral,
  repetições. Dê uma definição alternativa que garanta que o código
  de uma subárvore apareça apenas uma vez na lista resultante.
\end{prob}

\end{document}

